\documentclass[11pt, a4paper]{article}

% --- UNIVERSAL PREAMBLE BLOCK ---
% Geometry for A4 paper
\usepackage[a4paper, top=2.5cm, bottom=2.5cm, left=2cm, right=2cm]{geometry}

% Font and Language Setup
\usepackage{fontspec}
% Main language is english, secondary is chinese (for annotations)
\usepackage[english, bidi=basic, provide=*]{babel}

% Provide the languages
\babelprovide[import, onchar=ids fonts]{english}

% Set default/Latin font to Sans Serif (as per instructions)
% Noto Sans is a clean, modern font suitable for reports.
\setmainfont{Latin Modern Roman}
\setsansfont{Latin Modern Sans}
\setmonofont{Latin Modern Mono}

% --- DOCUMENT-SPECIFIC PACKAGES ---
\usepackage{amsmath}    % For advanced math environments
\usepackage{amssymb}    % For symbols like \mathbb
\usepackage{booktabs}   % For professional tables (\toprule, \midrule, \bottomrule)
\usepackage{graphicx}   % For images (though we use placeholders)
\usepackage{caption}    % For table/figure captions
\usepackage{fancyhdr}   % For headers and footers
\usepackage[
    colorlinks=true,
    linkcolor=blue,
    urlcolor=blue,
    citecolor=blue
]{hyperref} % For clickable links, MUST be loaded last

% --- PAGE STYLE ---
\pagestyle{fancy}
\fancyhf{} % Clear all header and footer fields
\fancyhead[L]{Huawei TechArena 2025: Phase II EMS Model}
\fancyhead[R]{Gen Li (Team SoloGen)}
\fancyfoot[C]{\thepage}
\renewcommand{\headrulewidth}{0.4pt}
\renewcommand{\footrulewidth}{0.4pt}

% --- TITLE ---
\title{Huawei TechArena 2025: BESS Energy Management System \\ \large Phase II Research Framework \& Mathematical Model}
\author{Gen Li (Team SoloGen) \\ \small \textit{Based on Phase I Model and Literature Review}}
\date{\today} % Or \date{October 31, 2025}

% --- DOCUMENT START ---
\begin{document}

\maketitle
\thispagestyle{fancy}

\begin{abstract}
This document outlines the mathematical formulation for the Phase II challenge of the Huawei TechArena 2025. This model directly extends the deterministic Mixed-Integer Linear Program (MILP) developed in Phase I. Responding to the practical constraints of the competition (deterministic data, open-source solvers), this framework explicitly avoids stochastic optimization.

The Phase II model introduces two key extensions:
\begin{enumerate}
    \item \textbf{Market Integration:} It incorporates the new \textbf{aFRR Energy Market}, co-optimizing energy bids alongside the Day-Ahead (DA) energy market and the FCR/aFRR capacity markets.
    \item \textbf{Degradation Modeling:} It replaces the simple "Daily Cycle Limit" (Cst-5) from Phase I with a sophisticated, linearized \textbf{degradation cost function}. This cost, subtracted from the objective, is based on the piecewise-linear models for both cyclic aging (Xu et al., 2017) and calendar aging (Collath et al., 2023), allowing the model to trade off revenue against battery lifetime dynamically.
\end{enumerate}
The resulting formulation remains a deterministic MILP, ensuring compatibility with open-source solvers while capturing the core economic and physical trade-offs of Phase II.
\end{abstract}

\section{Introduction}
The Phase I model established a strong foundation for co-optimizing BESS bids across three markets: DA Energy, FCR Capacity, and aFRR Capacity. The primary goal of Phase II is to enhance this model's realism and profitability by integrating two new factors: the aFRR Energy Market and the economic impact of battery degradation.

This document presents the revised mathematical formulation. It is designed to build directly upon the Phase I model, modifying existing constraints and adding new logic, rather than starting from a new, incompatible framework.

\section{Phase I Model - Base Model}

\subsection*{Objective Function}

The objective is to maximize the total net profit over the one-year horizon. This is the sum of day-ahead energy arbitrage revenue and ancillary service capacity payments, minus the cost of energy purchased for charging.

\begin{align}
\max\; Z
&= \sum_{t\in T} \Bigg( \frac{P_{DA}(t)}{1000}\, p_{\mathrm{dis}}(t) - \frac{P_{DA}(t)}{1000}\, p_{\mathrm{ch}}(t) \Bigg)\, \Delta t \notag \\
&\quad + \sum_{b\in B} \Big( P_{FCR}(b)\, c_{fcr}(b) + P^{\mathrm{pos}}_{aFRR}(b)\, c^{\mathrm{pos}}_{aFRR}(b) + P^{\mathrm{neg}}_{aFRR}(b)\, c^{\mathrm{neg}}_{aFRR}(b) \Big)\, \Delta b
\end{align}

\begin{itemize}
  \item The first term is day-ahead net profit. $p_{\mathrm{dis}}(t)$ and $p_{\mathrm{ch}}(t)$ are in kW, $P_{DA}(t)$ is in EUR/MWh; division by 1000 converts kW to MW. Multiplying by $\Delta t$ (0.25 h) yields EUR per interval.
  \item The second term is ancillary service capacity revenue. Bids $c(b)$ are in MW and prices $P(b)$ are in EUR/MW for each 4-hour block; multiplying by $\Delta b$ (4 h) yields EUR per block.
\end{itemize}

\subsection*{Constraint Equations with Detailed Explanations}

The optimization is subject to the following constraints.

\begin{enumerate}
  \item \textbf{State of Charge (SOC) Dynamics:} Energy balance of the BESS.
  \begin{align}
    e_{\mathrm{soc}}(t) &= e_{\mathrm{soc}}(t-1) + \Big( p_{\mathrm{ch}}(t)\,\eta_{\mathrm{ch}} - \frac{p_{\mathrm{dis}}(t)}{\eta_{\mathrm{dis}}} \Big)\, \Delta t
    \qquad \forall t\in T
  \end{align}
  For $t=1$, use the initial SOC:
  \begin{align}
    e_{\mathrm{soc}}(1) &= e^{\mathrm{init}}_{\mathrm{soc}} + \Big( p_{\mathrm{ch}}(1)\,\eta_{\mathrm{ch}} - \frac{p_{\mathrm{dis}}(1)}{\eta_{\mathrm{dis}}} \Big)\, \Delta t.
  \end{align}

  \item \textbf{SOC Limits:} Stay within operational boundaries.
  \begin{align}
    SOC_{\min}\,E_{\mathrm{nom}} \le e_{\mathrm{soc}}(t) \le SOC_{\max}\,E_{\mathrm{nom}} \qquad \forall t\in T
  \end{align}

  \item \textbf{Power Limits:} Binary–continuous linking to configured power.
  \begin{align}
    p_{\mathrm{ch}}(t) &\le y_{\mathrm{ch}}(t)\, P^{\mathrm{config}}_{\max} \qquad \forall t\in T \\
    p_{\mathrm{dis}}(t) &\le y_{\mathrm{dis}}(t)\, P^{\mathrm{config}}_{\max} \qquad \forall t\in T
  \end{align}

  \item \textbf{Simultaneous Operation Prevention:} No charge and discharge at the same time.
  \begin{align}
    y_{\mathrm{ch}}(t) + y_{\mathrm{dis}}(t) \le 1 \qquad \forall t\in T
  \end{align}

  \item \textbf{Market Co-optimization Power Limits:} Allocate total power between energy and reserves (convert MW bids to kW by $\times 1000$).
  \begin{align}
    p_{\mathrm{ch}}(t) + 1000\,c_{fcr}(b) + 1000\,c^{\mathrm{pos}}_{aFRR}(b) &\le P^{\mathrm{config}}_{\max} \qquad \forall b\in B,\, \forall t\in b \\
    p_{\mathrm{dis}}(t) + 1000\,c_{fcr}(b) + 1000\,c^{\mathrm{neg}}_{aFRR}(b) &\le P^{\mathrm{config}}_{\max} \qquad \forall b\in B,\, \forall t\in b
  \end{align}

  \item \textbf{Daily Cycle Limit:} Limit the total daily discharged energy.
  \begin{align}
    \sum_{t\in d} p_{\mathrm{dis}}(t)\,\Delta t \;\le\; N_{\mathrm{cycles}}\,E_{\mathrm{nom}}
    \qquad \forall d\in D
  \end{align}

  \item \textbf{Ancillary Service Energy Reserve:} Maintain sufficient energy to deliver awarded capacity for 15 minutes ($\Delta t=0.25$ h).
  \begin{align}
    e_{\mathrm{soc}}(t) &\ge SOC_{\min}\,E_{\mathrm{nom}} + \big(1000\,c_{fcr}(b) + 1000\,c^{\mathrm{pos}}_{aFRR}(b)\big)\,\Delta t
    && \forall b\in B,\, \forall t\in b \\
    e_{\mathrm{soc}}(t) &\le SOC_{\max}\,E_{\mathrm{nom}} - \big(1000\,c_{fcr}(b) + 1000\,c^{\mathrm{neg}}_{aFRR}(b)\big)\,\Delta t
    && \forall b\in B,\, \forall t\in b
  \end{align}

  \item \textbf{Minimum Bid Size Constraints:} Bids can be non-zero only if the corresponding binary is 1; if non-zero, they must satisfy minimum size and available power.
  \begin{align}
    c_{fcr}(b) &\ge y_{fcr}(b)\,\text{MinBid}_{fcr} && \forall b\in B \\
    c_{fcr}(b) &\le y_{fcr}(b)\,\frac{P^{\mathrm{config}}_{\max}}{1000} && \forall b\in B \\
    c^{\mathrm{pos}}_{aFRR}(b) &\ge y^{\mathrm{pos}}_{aFRR}(b)\,\text{MinBid}_{afrr} && \forall b\in B \\
    c^{\mathrm{pos}}_{aFRR}(b) &\le y^{\mathrm{pos}}_{aFRR}(b)\,\frac{P^{\mathrm{config}}_{\max}}{1000} && \forall b\in B \\
    c^{\mathrm{neg}}_{aFRR}(b) &\ge y^{\mathrm{neg}}_{aFRR}(b)\,\text{MinBid}_{afrr} && \forall b\in B \\
    c^{\mathrm{neg}}_{aFRR}(b) &\le y^{\mathrm{neg}}_{aFRR}(b)\,\frac{P^{\mathrm{config}}_{\max}}{1000} && \forall b\in B
  \end{align}

  \item \textbf{Ancillary Service Market Mutual Exclusivity:} Prevent simultaneous bidding in multiple ancillary service markets for the same block.
  \begin{align}
    y_{fcr}(b) + y^{\mathrm{pos}}_{aFRR}(b) + y^{\mathrm{neg}}_{aFRR}(b) \le 1 \qquad \forall b\in B
  \end{align}
\end{enumerate}

\section{Phase II Model Derivation}
We extend the Phase I MILP (Base Model) by modifying the objective function and relevant constraints.
% Set a TODO box here, say writing the cyclic and calendar aging models as model (ii) and (iii) step by step
\begin{itemize}
    \item[√] \textbf{Model (i):} Base Model + aFRR Energy Market
    \item[X] \textbf{Model (ii):} Model (i) + Cyclic Aging Cost
    \item[X] \textbf{Model (iii):} Model (ii) + Calendar Aging Cost = Full Phase II Model  
\end{itemize}

\subsection{Model (i): Base Model + aFRR Energy Market}
The aFRR Energy market (15-min granularity) operates similarly to the DA market as a price-taker energy market. We introduce new variables for bidding in this market and update all constraints that manage power and energy.

\begin{itemize}
    \item \textbf{New Variables:} $p^{\mathrm{pos}}_{aFRR,E}(t)$ (discharge/positive) and $p^{\mathrm{neg}}_{aFRR,E}(t)$ (charge/negative) for energy bids in the aFRR energy market.
    \item \textbf{New Binaries:} $y^{\mathrm{pos}}_{aFRR,E}(t)$ and $y^{\mathrm{neg}}_{aFRR,E}(t)$ to manage minimum bids.
    \item \textbf{Objective Function:} A new profit term $\mathbb{P}^{aFRR\_E}$ is added:
    $$ \mathbb{P}^{aFRR\_E} = \sum_{t\in T} \left( \frac{P^{\mathrm{pos}}_{aFRR, E}(t)}{1000}\, p^{\mathrm{pos}}_{aFRR, E}(t) - \frac{P^{\mathrm{neg}}_{aFRR, E}(t)}{1000}\, p^{\mathrm{neg}}_{aFRR, E}(t) \right)\, \Delta t $$
\end{itemize}

This requires revising the core physical constraints to prevent conflicts. We define new \textbf{total} charge/discharge variables:
$$ p^{\mathrm{total}}_{\mathrm{ch}}(t) = p_{\mathrm{ch}}(t) + p^{\mathrm{neg}}_{aFRR,E}(t) \quad \forall t $$
$$ p^{\mathrm{total}}_{\mathrm{dis}}(t) = p_{\mathrm{dis}}(t) + p^{\mathrm{pos}}_{aFRR,E}(t) \quad \forall t $$
These totals are then used to update the SOC, power, and mutual exclusivity constraints from Phase I.

\subsection{Integrating Battery Degradation as a Linear Cost}
We replace the rigid Phase I constraint (Cst-5) $\sum p_{\mathrm{dis}} \le N_{\mathrm{cycles}} E_{\mathrm{nom}}$ with a flexible and more profitable \textbf{degradation cost function}, $C^{\mathrm{Degradation}}$. This cost is subtracted from the objective function, allowing the optimizer to find the best balance between high-revenue, high-degradation actions and low-revenue, low-degradation actions.

Following the literature provided (Collath et al., 2023; Xu et al., 2017), we model this cost as the sum of calendar and cyclic aging:
$$ C^{\mathrm{Degradation}} = C^{\mathrm{cal}} + C^{\mathrm{cyc}} $$

\subsubsection{Cyclic Aging Cost ($C^{\mathrm{cyc}}$)}
Based on the methodology from Xu et al. (2017), we model cyclic aging as a \textbf{convex piecewise-linear cost of discharge}. The battery's energy capacity $E_{\mathrm{nom}}$ is divided into $J$ segments (e.g., 10 segments of 10\% SOC). Discharging from a shallower segment (e.g., 90-100\% SOC) is "cheaper" than discharging from a deeper segment (e.g., 20-30\% SOC).
\begin{itemize}
    \item \textbf{New Variables:} $e_{\mathrm{soc},j}(t)$ (energy in segment $j$), $p^{\mathrm{ch}}_{j}(t)$ (charge to segment $j$), and $p^{\mathrm{dis}}_{j}(t)$ (discharge from segment $j$).
    \item \textbf{Cost Function:} $C^{\mathrm{cyc}} = \sum_{t \in T} \sum_{j \in J} \left( c^{\mathrm{cost}}_{j} \cdot \frac{p^{\mathrm{dis}}_{j}(t)}{\eta_{\mathrm{dis}}} \cdot \Delta t \right)$
    \item \textbf{Logic:} The marginal cost $c^{\mathrm{cost}}_{j}$ increases as $j$ increases (deeper discharge), e.g., $c_1 < c_2 < ... < c_J$. The optimizer will naturally prefer to discharge from the "cheapest" (shallowest) available segment first, perfectly modeling the non-linear cost of deep discharges in a linear framework.
\end{itemize}

\subsubsection{Calendar Aging Cost ($C^{\mathrm{cal}}$)}
Based on Collath et al. (2023), calendar aging is a non-linear function of the average State of Charge (SOC). We linearize this cost using Special Ordered Sets of Type 2 (SOS2), a standard MILP technique.
\begin{itemize}
    \item \textbf{New Variables:} $\lambda_{t,i}$ (weighting variables for linearization).
    \item \textbf{Cost Function:} $C^{\mathrm{cal}} = \sum_{t \in T} c^{\mathrm{cal}}_{\mathrm{cost}}(t)$
    \item \textbf{Logic:} The total SOC $e_{\mathrm{soc}}(t)$ and its corresponding cost $c^{\mathrm{cal}}_{\mathrm{cost}}(t)$ are expressed as a convex combination of $I$ pre-calculated (SOC, Cost) breakpoints from the literature's non-linear curve.
    $$ e_{\mathrm{soc}}(t) = \sum_{i \in I} \lambda_{t,i} \cdot SOC^{\mathrm{point}}_i \quad | \quad c^{\mathrm{cal}}_{\mathrm{cost}}(t) = \sum_{i \in I} \lambda_{t,i} \cdot Cost^{\mathrm{point}}_i $$
    $$ \sum_{i \in I} \lambda_{t,i} = 1 \quad | \quad \{\lambda_{t,i}\} \text{ are SOS2 variables} $$
\end{itemize}
This formulation directly penalizes holding the battery at high SOC, as prioritized in the project description.

\section{Phase II Mathematical Formulation}

\subsubsection{Sets and Indices}
\begin{tabular}{@{}ll}
\toprule
$T, t$ & Set and index for 15-minute time intervals $t \in \{1, ..., 35040\}$ \\
$B, b$ & Set and index for 4-hour ancillary service blocks $b \in \{1, ..., 2190\}$ \\
$J, j$ & Set and index for cycle degradation cost segments $j \in \{1, ..., J\}$ \\
$I, i$ & Set and index for calendar degradation cost breakpoints $i \in \{1, ..., I\}$ \\
\bottomrule
\end{tabular}

\subsubsection{Parameters}
\begin{tabular}{@{}lll}
\toprule
Symbol & Description & Unit \\
\midrule
$P_{DA}(t)$ & Day-ahead electricity price & EUR/MWh \\
$P_{FCR}(b)$ & FCR capacity price & EUR/MW \\
$P^{\mathrm{pos}}_{aFRR}(b)$ & Positive aFRR capacity price & EUR/MW \\
$P^{\mathrm{neg}}_{aFRR}(b)$ & Negative aFRR capacity price & EUR/MW \\
$P^{\mathrm{pos}}_{aFRR,E}(t)$ & \textbf{(New)} Positive aFRR energy price & EUR/MWh \\
$P^{\mathrm{neg}}_{aFRR,E}(t)$ & \textbf{(New)} Negative aFRR energy price & EUR/MWh \\
$E_{\mathrm{nom}}$ & Nominal energy capacity & kWh \\
$P^{\mathrm{config}}_{\max}$ & Max charge/discharge power & kW \\
$\eta_{\mathrm{ch}}, \eta_{\mathrm{dis}}$ & Charging/discharging efficiencies & - \\
$SOC_{\min}, SOC_{\max}$ & Min/max SOC limits (fraction) & - \\
$E^{\mathrm{seg}}_{j}$ & \textbf{(New)} Max energy capacity of segment $j$ & kWh \\
$c^{\mathrm{cost}}_{j}$ & \textbf{(New)} Marginal cyclic degradation cost for segment $j$ & EUR/kWh \\
$SOC^{\mathrm{point}}_i$ & \textbf{(New)} SOC breakpoint $i$ for calendar cost & kWh \\
$Cost^{\mathrm{point}}_i$ & \textbf{(New)} Calendar degradation cost at breakpoint $i$ & EUR/hr \\
$\alpha$ & \textbf{(New)} Degradation price (meta-parameter) & - \\
\bottomrule
\end{tabular}

\subsubsection{Decision Variables}
\begin{tabular}{@{}lll}
\toprule
Symbol & Description & Type \\
\midrule
$p_{\mathrm{ch}}(t), p_{\mathrm{dis}}(t)$ & DA charge/discharge power & Cont. $\ge 0$ \\
$p^{\mathrm{pos}}_{aFRR,E}(t)$ & \textbf{(New)} aFRR-E positive (discharge) power & Cont. $\ge 0$ \\
$p^{\mathrm{neg}}_{aFRR,E}(t)$ & \textbf{(New)} aFRR-E negative (charge) power & Cont. $\ge 0$ \\
$c_{fcr}(b)$ & FCR capacity bid & Cont. $\ge 0$ \\
$c^{\mathrm{pos}}_{aFRR}(b)$ & Positive aFRR capacity bid & Cont. $\ge 0$ \\
$c^{\mathrm{neg}}_{aFRR}(b)$ & Negative aFRR capacity bid & Cont. $\ge 0$ \\
$p^{\mathrm{total}}_{\mathrm{ch}}(t), p^{\mathrm{total}}_{\mathrm{dis}}(t)$ & \textbf{(New)} Total charge/discharge power & Cont. $\ge 0$ \\
$p^{\mathrm{ch}}_{j}(t), p^{\mathrm{dis}}_{j}(t)$ & \textbf{(New)} Charge/discharge power for segment $j$ & Cont. $\ge 0$ \\
$e_{\mathrm{soc},j}(t)$ & \textbf{(New)} Energy stored in segment $j$ & Cont. $\ge 0$ \\
$e_{\mathrm{soc}}(t)$ & \textbf{(New)} Total energy stored in BESS & Cont. $\ge 0$ \\
$c^{\mathrm{cal}}_{\mathrm{cost}}(t)$ & \textbf{(New)} Calendar cost at time $t$ & Cont. $\ge 0$ \\
$\lambda_{t,i}$ & \textbf{(New)} SOS2 variable for calendar cost & Cont. $\ge 0$ \\
$y_{\mathrm{ch}}(t), y_{\mathrm{dis}}(t)$ & DA bid binaries & Binary \\
$y_{fcr}(b), y^{\mathrm{pos}}_{aFRR}(b), ...$ & Ancillary service bid binaries & Binary \\
$y^{\mathrm{pos}}_{aFRR,E}(t), ...$ & \textbf{(New)} aFRR-E bid binaries & Binary \\
$y^{\mathrm{total}}_{\mathrm{ch}}(t), y^{\mathrm{total}}_{\mathrm{dis}}(t)$ & \textbf{(New)} Total operation binaries & Binary \\
\bottomrule
\end{tabular}

\subsection{Objective Function}
The objective is to maximize the total profit, defined as market revenues minus degradation costs.
\begin{equation}
\max \; Z = \mathbb{P}^{DA} + \mathbb{P}^{ANCI} + \mathbb{P}^{aFRR\_E} - \alpha \cdot (C^{\mathrm{cyc}} + C^{\mathrm{cal}})
\end{equation}
Where:
\begin{align}
\mathbb{P}^{DA} &= \sum_{t\in T} \left( \frac{P_{DA}(t)}{1000}\, p_{\mathrm{dis}}(t) - \frac{P_{DA}(t)}{1000}\, p_{\mathrm{ch}}(t) \right)\, \Delta t \\
\mathbb{P}^{ANCI} &= \sum_{b\in B} \Big( P_{FCR}(b)\, c_{fcr}(b) + P^{\mathrm{pos}}_{aFRR}(b)\, c^{\mathrm{pos}}_{aFRR}(b) + P^{\mathrm{neg}}_{aFRR}(b)\, c^{\mathrm{neg}}_{aFRR}(b) \Big) \\
\mathbb{P}^{aFRR\_E} &= \sum_{t\in T} \left( \frac{P^{\mathrm{pos}}_{aFRR, E}(t)}{1000}\, p^{\mathrm{pos}}_{aFRR, E}(t) - \frac{P^{\mathrm{neg}}_{aFRR, E}(t)}{1000}\, p^{\mathrm{neg}}_{aFRR, E}(t) \right)\, \Delta t \\
C^{\mathrm{cyc}} &= \sum_{t \in T} \sum_{j \in J} \left( c^{\mathrm{cost}}_{j} \cdot \frac{p^{\mathrm{dis}}_{j}(t)}{\eta_{\mathrm{dis}}} \cdot \Delta t \right) \\
C^{\mathrm{cal}} &= \sum_{t \in T} c^{\mathrm{cal}}_{\mathrm{cost}}(t) \cdot \Delta t
\end{align}

\subsection{Constraints}

\subsubsection{Degradation-Aware SOC Dynamics (Replaces Cst-1)}
\begin{align}
e_{\mathrm{soc},j}(t) &= e_{\mathrm{soc},j}(t-1) + \left( p^{\mathrm{ch}}_{j}(t) \eta_{\mathrm{ch}} - \frac{p^{\mathrm{dis}}_{j}(t)}{\eta_{\mathrm{dis}}} \right) \Delta t & \forall t, \forall j \\
e_{\mathrm{soc}}(t) &= \sum_{j \in J} e_{\mathrm{soc},j}(t) & \forall t \\
p^{\mathrm{total}}_{\mathrm{ch}}(t) &= \sum_{j \in J} p^{\mathrm{ch}}_{j}(t) & \forall t \\
p^{\mathrm{total}}_{\mathrm{dis}}(t) &= \sum_{j \in J} p^{\mathrm{dis}}_{j}(t) & \forall t \\
p^{\mathrm{total}}_{\mathrm{ch}}(t) &= p_{\mathrm{ch}}(t) + p^{\mathrm{neg}}_{aFRR,E}(t) & \forall t \\
p^{\mathrm{total}}_{\mathrm{dis}}(t) &= p_{\mathrm{dis}}(t) + p^{\mathrm{pos}}_{aFRR,E}(t) & \forall t
\end{align}

\subsubsection{SOC \& Segment Limits (Replaces Cst-2)}
\begin{align}
SOC_{\min}\,E_{\mathrm{nom}} &\le e_{\mathrm{soc}}(t) \le SOC_{\max}\,E_{\mathrm{nom}} & \forall t \\
0 &\le e_{\mathrm{soc},j}(t) \le E^{\mathrm{seg}}_{j} & \forall t, \forall j
\end{align}

\subsubsection{Simultaneous Operation Prevention (Replaces Cst-3)}
\begin{align}
p^{\mathrm{total}}_{\mathrm{ch}}(t) &\le y^{\mathrm{total}}_{\mathrm{ch}}(t) \cdot P^{\mathrm{config}}_{\max} & \forall t \\
p^{\mathrm{total}}_{\mathrm{dis}}(t) &\le y^{\mathrm{total}}_{\mathrm{dis}}(t) \cdot P^{\mathrm{config}}_{\max} & \forall t \\
y^{\mathrm{total}}_{\mathrm{ch}}(t) + y^{\mathrm{total}}_{\mathrm{dis}}(t) &\le 1 & \forall t
\end{align}

\subsubsection{Market Co-optimization Power Limits (Replaces Cst-4)}
\begin{align}
p^{\mathrm{total}}_{\mathrm{dis}}(t) + 1000\,c_{fcr}(b) + 1000\,c^{\mathrm{pos}}_{aFRR}(b) &\le P^{\mathrm{config}}_{\max} & \forall b, t \in b \\
p^{\mathrm{total}}_{\mathrm{ch}}(t) + 1000\,c_{fcr}(b) + 1000\,c^{\mathrm{neg}}_{aFRR}(b) &\le P^{\mathrm{config}}_{\max} & \forall b, t \in b
\end{align}

\subsubsection{Ancillary Service Energy Reserve (Keeps Cst-6 Logic)}
\begin{align}
    \frac{\big(1000\,c_{fcr}(b) + 1000\,c^{\mathrm{pos}}_{aFRR}(b)\big)\,\tau}{\eta_{\mathrm{dis}}} &\leq e_{\mathrm{soc}}(t) - SOC_{\min}\,E_{\mathrm{nom}} & \forall b, t \in b \\
    \big(1000\,c_{fcr}(b) + 1000\,c^{\mathrm{neg}}_{aFRR}(b)\big)\,\tau\,\eta_{\mathrm{ch}} &\leq SOC_{\max}\,E_{\mathrm{nom}} - e_{\mathrm{soc}}(t) & \forall b, t \in b
\end{align}
\textit{(Where $\tau$ is the reserve duration, e.g., 0.25h)}

\subsubsection{Ancillary Capacity Market Exclusivity (Keeps Cst-7)}
\begin{equation}
    y_{fcr}(b) + y^{\mathrm{pos}}_{aFRR}(b) + y^{\mathrm{neg}}_{aFRR}(b) \le 1 \qquad \forall b
\end{equation}

\subsubsection{Cross-Market Mutual Exclusivity (Replaces Cst-8)}
\begin{align}
    y^{\mathrm{total}}_{\mathrm{ch}}(t) + y_{fcr}(b) + y^{\mathrm{pos}}_{aFRR}(b) &\le 1 & \forall b, t \in b \\
    y^{\mathrm{total}}_{\mathrm{dis}}(t) + y_{fcr}(b) + y^{\mathrm{neg}}_{aFRR}(b) &\le 1 & \forall b, t \in b
\end{align}

\subsubsection{Minimum Bids \& Binary Logic (Updates Cst-9)}
\begin{align}
    y_{\mathrm{ch}}(t)\,\text{MinBid}_{da} \cdot 1000 &\le p_{\mathrm{ch}}(t) \le y_{\mathrm{ch}}(t)\,P^{\mathrm{config}}_{\max} & \forall t \\
    y_{\mathrm{dis}}(t)\,\text{MinBid}_{da} \cdot 1000 &\le p_{\mathrm{dis}}(t) \le y_{\mathrm{dis}}(t)\,P^{\mathrm{config}}_{\max} & \forall t \\
    y^{\mathrm{pos}}_{aFRR,E}(t)\,\text{MinBid}_{afrr\_e} \cdot 1000 &\le p^{\mathrm{pos}}_{aFRR,E}(t) \le y^{\mathrm{pos}}_{aFRR,E}(t)\,P^{\mathrm{config}}_{\max} & \forall t \\
    y^{\mathrm{neg}}_{aFRR,E}(t)\,\text{MinBid}_{afrr\_e} \cdot 1000 &\le p^{\mathrm{neg}}_{aFRR,E}(t) \le y^{\mathrm{neg}}_{aFRR,E}(t)\,P^{\mathrm{config}}_{\max} & \forall t \\
    y^{\mathrm{total}}_{\mathrm{ch}}(t) &\ge y_{\mathrm{ch}}(t); \quad y^{\mathrm{total}}_{\mathrm{ch}}(t) \ge y^{\mathrm{neg}}_{aFRR,E}(t) & \forall t \\
    y^{\mathrm{total}}_{\mathrm{dis}}(t) &\ge y_{\mathrm{dis}}(t); \quad y^{\mathrm{total}}_{\mathrm{dis}}(t) \ge y^{\mathrm{pos}}_{aFRR,E}(t) & \forall t
\end{align}
\textit{(Similar min/max bid constraints for $c_{fcr}$, $c^{\mathrm{pos}}_{aFRR}$, $c^{\mathrm{neg}}_{aFRR}$ remain from Phase I)}

\subsubsection{Calendar Aging PWL Constraints (New)}
\begin{align}
    e_{\mathrm{soc}}(t) &= \sum_{i \in I} \lambda_{t,i} \cdot SOC^{\mathrm{point}}_i & \forall t \\
    c^{\mathrm{cal}}_{\mathrm{cost}}(t) &= \sum_{i \in I} \lambda_{t,i} \cdot Cost^{\mathrm{point}}_i & \forall t \\
    \sum_{i \in I} \lambda_{t,i} &= 1 & \forall t \\
    \{\lambda_{t,i}\}_{i \in I} &\text{ are SOS2 variables} & \forall t
\end{align}

\section{Implementation Strategy: Meta-Optimization \& Rolling Horizon}
The mathematical formulation above defines the "inner-loop" optimization. However, two external challenges must be solved to meet the competition's goals: \textbf{(1)} computational infeasibility of a full-year MILP solve, and \textbf{(2)} the need to select the "correct" trade-off between profit and degradation.

We address both with a "meta-optimization" strategy, which defines *how* we use the MILP model.

\subsection{Addressing Computational Feasibility: Rolling Horizon (MPC)}
The full-year model (T=35,040) is computationally intractable for open-source solvers. We will implement a \textbf{Rolling Horizon (or Model Predictive Control, MPC)} approach, as demonstrated in the literature (Collath et al., 2023).

The implementation will be a simulation loop:
\begin{enumerate}
    \item \textbf{Initialize:} Set $t=1$, $e_{\mathrm{soc},j}(0) = e^{\mathrm{init}}_{\mathrm{soc},j}$.
    \item \textbf{Loop (for $d = 1$ to 365):}
    \begin{itemize}
        \item \textbf{Solve:} Run the MILP defined in Section 3 for a short horizon, $T_{horizon}$ (e.g., 24-48 hours), starting from the current SOC $e_{\mathrm{soc},j}(t-1)$.
        \item \textbf{Execute:} Record the optimal bids and actions $p(t), c(b)$ for the first "execution" block (e.g., first 1 or 4 hours).
        \item \textbf{Update:} Calculate the new $e_{\mathrm{soc},j}$ at the end of the execution block. Set $t$ to the next time step.
    \end{itemize}
    \item \textbf{Aggregate:} Sum the profits and degradation from all 365 solves to get the total annual result.
\end{enumerate}
This makes the problem computationally feasible.

\subsection{Addressing the Bi-Objective Problem: The $\alpha$ Meta-Parameter}
The competition requires balancing two objectives: Revenue (30\% weight) and Degradation (30\% weight). Our objective function, $\max \; \mathbb{P} - \alpha \cdot C^{\mathrm{Degradation}}$, is a standard \textbf{Weighted Sum} method for solving this bi-objective problem.

The parameter $\alpha$ (the "degradation price") is the tuning knob that traces the Pareto frontier.
\begin{itemize}
    \item If $\alpha=0$, we only maximize profit (short life, high revenue).
    \item If $\alpha \to \infty$, we only minimize degradation (long life, low revenue).
\end{itemize}
The "correct" $\alpha$ is the one that results in the highest 10-Year ROI, as required by the "Investment Optimization" task. We will find this optimal $\alpha$ using a \textbf{meta-optimization loop} (a simple parameter sweep).

\subsubsection{Meta-Optimization Algorithm:}
\begin{enumerate}
    \item \textbf{Outer-Loop (Parameter Sweep):} For $\alpha$ in [0.1, 0.2, 0.3, ..., 3.0]:
    \begin{itemize}
        \item \textbf{Inner-Loop (Simulation):} Run the full 365-day \textbf{Rolling Horizon (MPC)} simulation using this value of $\alpha$.
        \item \textbf{Record:} Store the resulting total annual profit $\mathbb{P}(\alpha)$ and total annual SOH loss $\Delta SOH(\alpha)$.
        \item \textbf{Project:} Calculate the 10-Year ROI for this $\alpha$, factoring in the SOH loss (which determines battery replacement, if any).
    \end{itemize}
    \item \textbf{Select:} Choose the $\alpha^*$ that produced the \textbf{maximum 10-Year ROI}.
    \item \textbf{Final Run:} Run the simulation one last time with $\alpha^*$ to generate the final bidding files.
\end{enumerate}
This meta-optimization strategy directly links our MILP model to the competition's final evaluation criteria (10-year ROI) and correctly implements the bi-objective optimization you identified.

\section{Conclusion and Next Steps}
This Phase II formulation successfully integrates the aFRR energy market and a sophisticated, linearized degradation cost model into the Phase I MILP. The implementation strategy (MPC) ensures computational feasibility, while the meta-optimization of $\alpha$ ensures we find the optimal balance between profit and degradation to maximize the 10-year ROI.

\begin{thebibliography}{9}
\bibitem{xu2017}
B. Xu, J. Zhao, T. Zheng, E. Litvinov, and D. S. Kirschen. (2017). "Factoring the Cycle Aging Cost of Batteries Participating in Electricity Markets." \textit{arXiv:1707.04567v2}.

\bibitem{collath2023}
N. Collath, M. Cornejo, V. Engwerth, H. Hesse, and A. Jossen. (2023). "Increasing the lifetime profitability of battery energy storage systems through aging aware operation." \textit{Applied Energy}, 348, 121531.

\bibitem{huawei2025}
Gen Li. (2025). "Huawei TechArena 2025: BESS Energy Management System." \textit{Internal Project Document}.
\end{thebibliography}

\end{document}

